% Section 4 draft: Optimization Strategy. Paste over the current \section{Optimization Strategy} block.
% References: docs/paper/section4_references.bib. EDIT: marks facts to check before submission.

\section{Optimization Strategy}
\vspace{0.35em}

The optimizations were developed on a small, widely used detector and then
applied to the stakeholder's SAR detector. This order was first a practical
one: the stakeholder's model and data became available only well into the
project, so the replica, the measurement procedure and the optimizations were
built on a public model in the meantime. It also separates two questions:
whether a change removes the bottleneck it targets, and whether the gain
carries over to a heavier model with a different balance between GPU work and
pipeline overhead.

\subsection{Detection Models}

\textbf{YOLOv8n} is the smallest model of the Ultralytics YOLOv8
family~\cite{ultralytics-yolov8}. We use the COCO-pretrained weights
unchanged: 80 classes, a three-channel 640$\times$640 input, and detection
heads at strides 8, 16 and 32, which give 8400 candidate boxes per frame. It
served as the development model while the stakeholder's model was not yet
available: it is light, well documented and runs through the standard
Ultralytics inference code.

\textbf{M15} is the stakeholder's detector for airborne SAR imagery. It is
based on YOLO11s~\cite{ultralytics-yolo11}, adapted to one class (person) and
a single-channel grayscale input of 960$\times$960, and extended with an
additional detection head at stride 4 (P2) for small targets. With four
heads, it produces 76\,500 candidate boxes per frame. We received a training
checkpoint, rebuilt the architecture as an Ultralytics model, loaded the
weights with a strict check that every tensor matches, and exported it
through ONNX to TensorRT. Table~\ref{tab:models} compares the two models.

\begin{table}[H]
  \centering
  \small
  \begin{tabular}{@{}lll@{}}
    \toprule
    & YOLOv8n & M15 \\
    \midrule
    Base architecture & YOLOv8n & YOLO11s + P2 head \\
    Input & 640$\times$640, 3 channels & 960$\times$960, 1 channel \\
    Classes & 80 (COCO) & 1 (person) \\
    Detection strides & 8, 16, 32 & 4, 8, 16, 32 \\
    Candidates per frame & 8\,400 & 76\,500 \\
    Weights & COCO-pretrained & Stakeholder's checkpoint \\
    \bottomrule
  \end{tabular}
  \vspace{0.5em}
  \caption{The two detection models.}
  \label{tab:models}
\end{table}
\vspace{-0.6em}

\subsection{Baseline Inference Pipeline}

Both baselines run a TensorRT engine built on the Jetson from the ONNX
export, in FP32 with TensorRT allowed to use TF32 for matrix operations. Both
read the 1920$\times$1080 video with OpenCV on the CPU and use a confidence
threshold of 0.25 and an NMS IoU threshold of 0.7, the Ultralytics default.

The YOLOv8n baseline is the unmodified Ultralytics inference path with the
TensorRT engine: preprocessing (letterboxing, color conversion and
normalization) on the CPU, one engine execution per frame in which every
kernel is launched individually, and Ultralytics' NMS on the GPU. The M15
baseline follows the same structure in our own script, because the
checkpoint is not an Ultralytics model file: OpenCV converts the frame to
grayscale, resizes it and pads it to 960$\times$960 on the CPU; the tensor is
copied into the engine's input buffer; the engine runs; and
\texttt{torchvision} performs class-agnostic NMS. Every later stage of M15
reuses these functions, so a change in throughput can be attributed to the
one step that changed.

\subsection{Profiling-Guided Optimization Stages}

Each stage followed the same loop: profile the current pipeline with Nsight
Systems, find the largest avoidable cost on the timeline (idle gaps between
kernels, CPU work in front of the GPU, or the kernels themselves), change
that one step, and measure again. The stages are cumulative: each keeps the
changes before it (Table~\ref{tab:optimization-stages}).

\textbf{CUDA graph.} In the baseline, the CPU launches every kernel of the
network separately on each frame. When the kernels are short, the GPU
finishes one before the CPU has issued the next, which leaves idle gaps on
the timeline. A CUDA graph records the sequence of kernel launches once and
replays it with a single call. We capture the TensorRT engine's execution
into a graph after one warm-up run and replay it on every frame. A captured
graph reads from and writes to the same memory addresses on every replay and
assumes the same tensor shapes, and operations that make the CPU wait for the
GPU are not allowed during capture~\cite{pytorch-cuda-graphs}. These rules
decide the design of the next two stages: GPU preprocessing writes into the
engine's fixed input buffer, and NMS has to be rewritten with fixed shapes
before it can be captured.

\textbf{GPU preprocessing.} After the graph, preprocessing on the CPU became
the largest cost in front of the network. In this stage the decoded 8-bit
frame is uploaded once, and resizing, color conversion, normalization and
letterbox padding run on the GPU, writing straight into the engine's input
buffer. The order of operations matters: converting or normalizing the full
1920$\times$1080 frame before resizing processes several times more pixels
than necessary. Resizing first and converting the smaller result reduced the
GPU time of M15's preprocessing from 5.4--7~ms to 1.3~ms per frame.
% EDIT: the 5.4-7 ms -> 1.3 ms figure is from the profiled runs (docs/m15.md); keep or move to Section 6.

\textbf{Fixed-shape NMS.} Standard NMS first keeps the boxes above the
confidence threshold, so the size of every following tensor depends on the
image. Each such step makes the CPU wait for the GPU to learn the size before
it can continue, and the variable shapes prevent capturing NMS into a graph.
The postprocessing time also grows with the number of detections, so frames
with targets take longer. We replace it with a fixed-shape version: the
$k$ highest-scoring candidates are selected (\texttt{topk}),
suppression is computed as a $k \times k$ IoU matrix, and the output is
always a $k \times 6$ tensor in which unused rows have a score of zero.
Nothing in it depends on the data, so the CPU can queue all of its operations
without waiting, and the only synchronization is the single copy of the
result to the CPU at the end. This removes the waits even when NMS runs
outside the graph; capturing it as well also removes the cost of launching
its operations one by one. The fixed version costs the same on every frame.
Its gain therefore grows with the number of candidates above the threshold,
up to $k$, while on frames without detections, where standard NMS stops
early, it can cost slightly more. Its main effect is on the slowest frames
rather than on the mean, and the size of the gain depends on how many
targets the video contains.

Both models use the same suppression rule, Cluster-NMS~\cite{zheng2020clusternms}:
a box is removed only if a higher-scoring box that is itself kept overlaps it,
repeated a fixed number of times. This gives exactly the result of standard
NMS while keeping the shapes fixed. For YOLOv8n's 80 classes, boxes are
shifted by their class so that boxes of different classes never suppress each
other, and $k = 300$, Ultralytics' limit on detections per image; both models
use the baseline's IoU threshold of 0.7. Our first version for YOLOv8n used
Fast NMS~\cite{bolya2019yolact}, a single matrix step that also lets removed
boxes suppress others, with $k = 100$; on COCO it cost about four points of
mAP50, mostly on small objects, which is why both models now use the exact
rule. For M15 ($k = 100$) we go one step further and capture the NMS into the
same CUDA graph as the engine, so a frame's network and postprocessing replay
with one call; for YOLOv8n it runs after the graph replay. In our test runs no
M15 frame had more than $k$ candidates above the threshold, so its fixed NMS
matched standard NMS on every frame.

\textbf{FP16.} The last stage builds the engine with 16-bit floating-point
layers. Ultralytics' FP16 export keeps FP32 input and output tensors, so
TensorRT adds conversion kernels at both ends of the network on every frame.
We therefore build the engines with FP16 input and output as well: the
preprocessing writes FP16 directly into the input buffer, and the network
output stays in FP16 through candidate selection. One conversion to FP32 is
kept on purpose, after the top-$k$ selection: box areas in the
960$\times$960 input reach 921\,600, beyond FP16's largest value of 65\,504,
and NMS computes areas and overlaps in the tensor's precision. Because FP16
changes the numbers the network computes, its effect on the detections is
measured on the SAR test set (Section~6).

\begin{table}[H]
  \centering
  \small
  \begin{tabular}{@{}p{0.31\columnwidth}p{0.63\columnwidth}@{}}
    \toprule
    Stage & Targeted bottleneck \\
    \midrule
    TensorRT baseline & Reference configuration for the staged comparison \\
    + CUDA graph & CPU launch overhead and inter-kernel gaps \\
    + GPU preprocessing & CPU preprocessing and associated data movement \\
    + Fixed-shape NMS & Variable postprocessing work and GPU--CPU synchronization \\
    + FP16 & GPU kernel execution time and data movement \\
    \bottomrule
  \end{tabular}
  \vspace{0.5em}
  \caption{Cumulative optimization stages and their targeted bottlenecks.}
  \label{tab:optimization-stages}
\end{table}
\vspace{-0.6em}
